
The comparative analysis and subsequent mechanistic investigation of the simulation results yielded a definitive conclusion regarding the optimization of solid waste management policies in the Municipality of Bacolod. The primary finding was that the conventional ``Status Quo'' approach, characterized by an equitable distribution of resources across all barangays with a heavy emphasis on Information, Education, and Communication (IEC), was mathematically insufficient to achieve municipality-wide compliance. While this strategy was capable of maintaining order in smaller, socially cohesive communities, it consistently failed to overcome the high behavioral friction present in densely populated urban centers. This failure was not due to a lack of total funding, but rather a strategic error in resource dilution; spreading the budget thinly prevented the Local Government Unit (LGU) from ever reaching the critical enforcement intensity required to break the inertia of non-compliance in high-resistance zones.

Furthermore, the simulation demonstrated that single-instrument policies---whether purely punitive or purely incentive-based---were equally flawed when applied uniformly. Pure enforcement collapsed due to its inability to sustain positive attitudes, while pure incentives failed due to the diminishing per-capita value of rewards in large populations. In contrast, the Deep Reinforcement Learning (DRL) agent's ``Sequential Saturation'' strategy provided a robust solution to this dilemma. By abandoning the heuristic of simultaneous equality in favor of sequential efficacy, the agent successfully leveraged the ``Graduation Effect,'' utilizing social norms to sustain compliance in stabilized areas while concentrating capital to besiege and conquer the most resistant urban targets. This confirmed that the optimal policy for a resource-constrained LGU was dynamic and adaptive: a phased intervention that treated compliance not as a static maintenance task, but as a series of strategic tipping points to be systematically triggered.